\documentclass[10pt,a4paper]{article}

\usepackage[a4paper,top=0.92cm,bottom=0.92cm,left=1.28cm,right=1.18cm]{geometry}
\usepackage{iftex}
\ifPDFTeX
  \usepackage[utf8]{inputenc}
  \usepackage[T1]{fontenc}
  \usepackage{CJKutf8}
\else
  \usepackage[UTF8]{ctex}
\fi

\usepackage{xcolor}
\usepackage{tabularx}
\usepackage{array}
\usepackage{enumitem}
\usepackage{graphicx}
\usepackage[hidelinks]{hyperref}

\setlength{\parindent}{0pt}
\setlength{\parskip}{2.0pt}
\linespread{1.14}
\setlist[itemize]{leftmargin=1.5em,label={$\bullet$},itemsep=1.6pt,topsep=1.1pt,parsep=0pt,partopsep=0pt}
\pagestyle{empty}

\definecolor{SectionGray}{RGB}{249,250,252}
\definecolor{TextGray}{RGB}{90,90,90}
\definecolor{MainBlueGray}{RGB}{58,72,88}
\definecolor{SectionTitleBlue}{RGB}{46,78,112}

\newcommand{\sectionbar}[1]{%
  \vspace{3.2pt}
  \begingroup
  \setlength{\fboxsep}{1.1pt}%
  \noindent\colorbox{SectionGray}{\parbox{\dimexpr\linewidth-2\fboxsep\relax}{\textbf{\textcolor{SectionTitleBlue}{#1}}}}%
  \endgroup
  \vspace{1.8pt}
}

\newcommand{\majorsection}[1]{%
  \vspace{0.62\baselineskip}
  \sectionbar{#1}
  \vspace{0.42\baselineskip}
}

\newcommand{\entry}[3]{%
  \noindent\begin{tabularx}{\linewidth}{@{}>{\raggedright\arraybackslash}p{0.24\linewidth}>{\centering\arraybackslash}p{0.48\linewidth}>{\raggedleft\arraybackslash}p{0.28\linewidth}@{}}
  {\textbf{\fontsize{10.8}{12}\selectfont #1}} & \textbf{#2} & {\textbf{\fontsize{10.4}{12}\selectfont #3}}
  \end{tabularx}\\[-1pt]
}

\newcommand{\entryedu}[3]{%
  \noindent\begin{tabularx}{\linewidth}{@{}>{\raggedright\arraybackslash}p{0.24\linewidth}>{\centering\arraybackslash}p{0.48\linewidth}>{\raggedleft\arraybackslash}p{0.28\linewidth}@{}}
  {\textbf{\fontsize{10.8}{12}\selectfont #1}} & \textbf{#2} & {\textbf{\fontsize{10.2}{12}\selectfont \mbox{#3}}}
  \end{tabularx}\\[-1pt]
}

\newcommand{\projdesc}[1]{%
  \noindent\textcolor{TextGray}{\textbf{项目描述：}}#1\\[1.1pt]
}

\newcommand{\projcore}{%
  \noindent\textcolor{TextGray}{\textbf{核心工作：}}\\[0.0pt]
}

\begin{document}
\ifPDFTeX
\begin{CJK*}{UTF8}{gbsn}
\fi

% ===== 顶部信息区 =====
\noindent\begin{minipage}[t]{0.75\linewidth}
  \vspace{0pt}
  {\LARGE \textbf{王世卓}}\\[6pt]
  \renewcommand{\arraystretch}{1.2}
  \begin{tabularx}{\linewidth}{@{}p{1.1cm}X p{1.1cm}X@{}}
    年\;龄 & 22岁 & 性\;别 & 男 \\
    电\;话 & 18837502016 & 邮\;箱 & 2151056619@qq.com \\
  \end{tabularx}\\[4pt]
  {\normalsize \textbf{求职意向：软件工程师（AI应用 / Agent方向）}}
\end{minipage}%
\hfill
\begin{minipage}[t]{0.20\linewidth}
  \vspace{0pt}
  \raggedleft
  \IfFileExists{avatar1.jpg}{%
    \includegraphics[width=2.8cm,keepaspectratio]{avatar1.jpg}
  }{%
    \fbox{\parbox[c][3.5cm][c]{2.8cm}{\centering 照片}}
  }
\end{minipage}
\vspace{0.2\baselineskip}

% ===== 教育背景 =====
\majorsection{教育背景}
\entryedu{2022-09 -- 2026-06}{西南石油大学（双一流）}{数据科学与大数据技术（本科）}
\noindent\textbf{主修课程：} \textcolor{TextGray}{操作系统、计算机网络、面向对象程序设计、数据结构、数据库系统}
\vspace{0.3\baselineskip}

% ===== 专业能力 =====
\majorsection{专业能力}
\begin{itemize}
  \item \textbf{赴日意愿与学习力}：职业规划明确，期望赴日长期发展。具备较强自驱力与问题拆解能力，能全力配合公司日语培训，按期达到业务沟通标准。
  \item \textbf{AI 应用与开发}：具备扎实 Python 基础，熟悉 ChatGPT / DeepSeek 等主流模型；熟练完成 Prompt 优化，具备Agent/RAG/Workflow 架构的实战能力。
  \item \textbf{工程技术栈}：熟悉FastAPI/Redis/Docker 等技术栈；熟悉 \textbf{Linux 与 Git}，能快速查阅英文文档并独立完成接口开发、日志定位与问题排查。
\end{itemize}

% ===== 项目经历 =====
\majorsection{项目经历}
\entry{2025-07 -- 2025-09}{多 Agent 智能业务工单处理系统}{个人项目}
\projdesc{针对售后工单处理场景，基于 Python + LangChain + FastAPI 搭建多 Agent 协作系统，实现用户文本请求的任务拆解、动作生成与结果校验。}
\projcore
\begin{itemize}
  \item 设计并实现规划 Agent、执行 Agent、质检 Agent 协作流程，分别完成任务拆解、动作生成与结果校验。
  \item 基于 JSON Schema 约束模型输出结构，针对非标准返回加入参数校验、自动重试与降级兜底机制，降低执行失败率。
  \item 针对多轮工单交互场景，实现对话摘要压缩与上下文裁剪，减少冗余 Token 消耗，提升处理稳定性。
  \item 支持工单分类、字段提取、动作生成等核心场景，完成基本工单处理流程。
\end{itemize}

\entry{2025-12 -- 2026-02}{基于 n8n Workflow 的多源信息采集与分析系统}{个人项目}
\projdesc{面向个人与小团队的信息跟踪需求，基于 n8n Workflow 自动采集 RSS/API 内容，完成摘要提取、结构化整理、推送与归档，减少重复信息筛选与同步成本。}
\projcore
\begin{itemize}
  \item 搭建信息采集与处理 Workflow，整合 RSS/API 数据源，通过 HTTP Request 与 Webhook 节点实现外部内容采集、清洗与标准化入库。
  \item 基于提示词模板设计结构化抽取流程，调用 DeepSeek API 输出主题、摘要、来源等字段，提升后续检索与复用效率。
  \item 通过自定义 JavaScript 节点对模型返回的非标准 JSON 进行二次修复与字段校验，减少推送失败和入库异常。
  \item 打通“采集—分析—推送—归档”自动化处理流程，提升信息整理效率，为持续跟踪行业与技术动态提供支持。
\end{itemize}

% ===== 自我评价 =====
\majorsection{自我评价}
\begin{itemize}[itemsep=1.2pt,topsep=0.9pt]
  \item \textbf{喜爱日本文化与生活环境}：具备明确且坚定的赴日意愿，向往当地的生活氛围，拥有较强的跨文化适应力与长期留任定力。
  \item \textbf{扎实的技术底盘}：计算机基础扎实，对 AI 应用抱有热情，具备较强的技术自驱力，习惯独立阅读源码与英文文档解决复杂落地问题。
  \item \textbf{闭环的工程思维}：做事踏实、责任心强，注重交付闭环。生活独立性高，抗压能力较强，能快速融入跨文化团队并提供稳定产出。
\end{itemize}

\ifPDFTeX
\end{CJK*}
\fi
\end{document}
